\input{estilo}
% Mejorada: los dos nodos uno debajo del otro, para que se vea que reciben lo
% mismo; los dos parten de S0 (adentro del cilindro), y los tres requisitos
% con una llave que los abarca.
\begin{document}
\begin{figura}\begin{tikzpicture}[celda/.append style={minimum height=7mm}]
  \foreach \y/\n/\k in {1.5/a/1, 0/b/2} {
    \tiracontorno[0.7]{\n}{(0,\y)}{\var{o}\textsubscript{1}, \var{o}\textsubscript{2},
      \var{o}\textsubscript{3}, \var{o}\textsubscript{4}}
    \draw[msg] (\n-4.east) -- ++(0.5,0);
    \pic[scale=0.85] at (3.75,{\y-0.12}) {base};
    \node[font=\small] at (3.75,{\y+0.3}) {\var{S}\textsubscript{0}};
    \node[font=\small, anchor=east] at (-0.2,{\y+0.35}) {nodo \k};
  }
  \draw[llave=6pt] (4.6,2.0) -- (4.6,0.2);
  \foreach \t [count=\i] in {operaciones en el mismo orden, determinismo estricto,
      {mismo estado inicial \var{S}\textsubscript{0}}}
  {\fill (5.0,{1.1+0.55*(2-\i)}) circle (1.6pt);
    \node[font=\small, anchor=west] at (5.15,{1.1+0.55*(2-\i)}) {\t};}
\end{tikzpicture}\end{figura}
\end{document}
